% Einheit 2 von 5 – Streckenlängen, Mittelpunkt und Teilpunkte (bank/punkte-und-strecken-im-koordinatensystem/e2.jsonl, zusammenbau v0.1)
%% TODO Zeitmarke Einheit 2: Marken-Zeile nicht lesbar
%% TODO Zweigzeile Teil 1: was hier gelernt wird (Satz in Schülersprache aus dem Einheitstitel) – Einheit 2
\einheitenkopf[e2]{Einheit 2 von 5 · Streckenlängen, Mittelpunkt und Teilpunkte}
\zweigzeile{Abitur GK}
\setcounter{aufgabe}{12}

%% TODO Titel: Ich-kann-Satz fehlt in der Bank (Platzhalter: Kettenname) – Nr. 13
%% TODO Anweisung: ein Satz über der ersten Teilaufgabe fehlt in der Bank – Nr. 13
\begin{aufgabe}{Streckenlänge, Mittelpunkt, Teilpunkte}
%% TODO \rechenplatz{n}: Schrittzahl des Grundfalls fehlt in der Bank (nur bei fester Schreibform, 2.3 b)
\begin{teile}
\teil Wo muss der Haken an der Stange zwischen $A$ und $B$ hängen, damit er genau in der Mitte ist? Was ist gefragt? \\ \kreuz{eine Länge} \\ \kreuz{ein Punkt} \\ \kreuz{eine Figur}
\teil $A(2 | 1 | 3)$, $B(4 | 4 | 9)$: Länge von $AB$? \leerfeld LE
\teil Ein Zeltmast endet in $S(2 | 2 | 6)$; vier gleich lange Leinen führen zu Heringen am Boden, einer davon steckt in $H(4 | 5 | 0)$ (1 LE = 1 m). Je Leine kommen $40$ cm für die Knoten dazu. Wie viel Leine braucht man insgesamt? \leerfeld[m]
\teil Ein Gummiseil ist von $A(3 | 1 | 2)$ nach $B(9 | 9 | 2)$ gespannt (1 LE = 1 m) und dabei um $25\,\%$ länger als ungespannt. Wie lang ist es ungespannt? \hfill (Abitur 2020 GK) \\ \leerfeld[m]
\teil Eine Pyramide ist symmetrisch zur $x_1x_3$-Ebene; $D(1 | 3 | 0)$ ist eine Grundecke, $S(3 | 0 | 6)$ die Spitze. Zeige: $W(2 | 1{,}5 | 3)$ ist der Mittelpunkt der Kante $DS$. Gib den Mittelpunkt $T$ der zu $DS$ symmetrischen Kante an. \hfill (Abitur 2026 GK) \\ T( \leerfeld | \leerfeld | \leerfeld )
\teil Eine Strebe führt von $P(9 | 0 | 3)$ nach $Q(0 | 3 | 0)$; zwei Schellen teilen sie in drei gleich lange Stücke. Koordinaten der Schellen? \hfill (Abitur 2024 GK)
\teil Eine Stütze führt von $M(3 | 3 | 0)$ zur Spitze $D(0 | 0 | 9)$. $Q_h$ ist der Punkt der Strecke in der Höhe $h$. Koordinaten von $Q_h$? Gib $Q_{3}$ an. \hfill (Abitur 2022 LK)
\teil $A(3 | 5 | 2)$ und $B(5 | 6 | 4)$ liegen auf einer Geraden $g$. Gib einen Punkt $D \ne B$ auf $g$ an, der von $A$ so weit entfernt ist wie $B$. \hfill (Abitur 2020 GK) \\ D( \leerfeld | \leerfeld | \leerfeld )
\teil $A(2 | -1 | 2)$; der Punkt $C$ ist festgelegt durch $\overrightarrow{OC} = 3 \cdot \overrightarrow{OA}$. Länge der Strecke $AC$? \hfill (Abitur 2018 GK) \\ \leerfeld
\teil Die Ecken eines Oktaeders sind die Mittelpunkte der Seitenflächen eines Würfels; $A(2 | 1 | 3)$ und $C(-2 | -3 | 5)$ sind gegenüberliegende Ecken des Oktaeders. Zeige: Die Würfelkante ist $6$ lang. \hfill (Abitur 2024 LK)
\end{teile}
\end{aufgabe}

%% TODO Titel: Ich-kann-Satz fehlt in der Bank (Platzhalter: Kettenname) – Nr. 14
%% TODO Anweisung: ein Satz über der ersten Teilaufgabe fehlt in der Bank – Nr. 14
\begin{aufgabe}{Streckenlänge, Mittelpunkt, Teilpunkte – Fehler finden, Begründen}
\begin{teile}
\teil Ich finde den Fehler: $A(4 | 1 | 3)$, $B(6 | 2 | 5)$. Jan rechnet $|\overrightarrow{AB}| = \sqrt{6^2 + 2^2 + 5^2} = \sqrt{65}$.
\teil Begründe: Die Länge von $\vec v = (a | b | c)$ ist $\sqrt{a^2 + b^2 + c^2}$ – das ist der Satz des Pythagoras im Raum.
\end{teile}
\end{aufgabe}
